Nguồn: \href{https://1thegioi.vn/giao-su-phan-dinh-dieu-nhu-toi-biet-15579.html}{Báo Một Thế Giới}

Tôi nói chuyện với ông Phan Đình Diệu khá nhiều lần tại Paris trong khoảng thời gian những năm 1989, 1990, 1991...

Chắc nhiều người còn nhớ khoảng thời gian năm 1989-1990 các nước Đông Âu rục rịch thay đổi để rồi từng nước, từng nước rời bỏ khối xã hội chủ nghĩa. Đây cũng là thời gian mà Đảng Cộng sản Việt Nam, dưới sự lãnh đạo của Tổng bí Thư Nguyễn Văn Linh, cố sức chèo chống để Việt Nam có thể trụ được cùng với các nước xã hội chủ nghĩa còn sót lại. Thời gian đó các trí thức Việt trong nước bắt đầu ra nước ngoài nhiều hơn.

Một lần lên Quán Cơm Việt Nam tại Place de Maubert-Mutualité. Quán đông, trong lúc nhìn quanh tìm chỗ, một người đàn ông thấp thấp đang dùng cơm chiều mời tôi: “Mời anh ngồi chung bàn!”. Cái cơ duyên cơm chiều đó khiến tôi được quen với một người mà tôi từng khâm phục khi nghe một số bậc trưởng thượng nhắc lại câu nói chục năm trước của ông về một nhà lãnh đạo Việt Nam đang lúc đầy quyền uy, đại ý “Đồng chí là nhà lãnh đạo giành độc lập dân tộc vĩ đại. Và còn vĩ đại hơn nếu lúc này đồng chí từ chức”.

Phan Đình Diệu là nhà toán học nổi tiếng. Vài người bạn đang dạy toán tại Trường đại học Orsay và Trường Paris 7 cho tôi biết “Tay đó giỏi lắm, Việt Nam cô lập vậy mà vẫn nắm được những kiến thức mới”. Tôi không hiểu nhiều về lĩnh vực chuyên môn của ông, và cũng không có ý định hỏi thêm vì nghĩ mình không hiểu. Phần ông, ông lại muốn thăm nơi tôi làm việc. Tôi đã đưa ông thăm phòng thí nghiệm của khoa học về cá (Laboratoire d’Ichthyologie) của Viện bảo tàng Lịch sử tự nhiên Paris, và phòng thí nghiệm Sinh học phân tử Trường đại học Orsay. Ông chăm chú lắng nghe, và rất thích thú với lịch sử tiến hóa của sinh vật.

Tại góc uống cà phê của phòng thí nghiệm, ông hỏi: “Thế giới vật chất quanh ta bao gồm giới sống và giới không sống (living things and non-living things). Tại sao giới sống phát triển muôn hình vạn trạng, phát sinh hàng chục triệu loài sinh vật trong khi giới không sống thì cứ trơ ra?”. Xin mở ngoặc: giới sống gồm tất cả các dạng sinh vật, còn giới không sống gồm tất cả các vật còn lại như đất, đá, nước, các loại khoáng sản...

Câu hỏi ông nêu ra cũng là một câu hỏi căn bản của lịch sử tiến hóa sinh giới. Chúng tôi thảo luận, đúng hơn là tôi trình bày cho nghe ông những kiến thức mới nhất trong ngành mà tôi biết cho tới lúc đó.

Trước khi ra căng tin dùng buổi trưa, ông gật gù: “Sinh vật thích nghi và phát triển vô số loài vì chúng mềm dẻo, cộng tác nhau. Chúng có nhiều mối quan hệ chấp nhận chứ không loại bỏ nhau: cộng sinh, ký sinh, hỗ sinh... Nhiều mối quan hệ chấp nhận nhau tạo sự bền vững và mềm dẻo để tồn tại, thích nghi và tiến hóa. Sinh vật tiến hóa được vì khác với đất đá là ở điểm này”.

Tôi hiểu ngay đề tài được ông học hỏi từ lâu, và ông muốn đồng quy các kiến thức về khoa học tự nhiên với khoa học xã hội. Quả thực ông có tầm vóc một nhà trí thức lớn.

Với cái nhìn thuận chiều tiến hóa như vậy, Phan Đình Diệu không khiến tôi ngạc nhiên khi ông theo xu hướng đổi mới chính trị mạnh mẽ. Khi thấy các nước Ba Lan, Tiệp Khắc, Hung... chuyển đổi mạnh mẽ, quyết liệt, ông thể hiện rõ ràng quan điểm ủng hộ: “Anh xem, họ sẽ giàu mạnh rất nhanh!”. Và ông tiếc: “Việt Nam mình mấy năm trước nổi tiếng cởi mở nhất trong các nước xã hội chủ nghĩa. Lúc đó mình sang Tiệp, Hung, các bạn thèm thuồng nhìn Việt Nam như hình mẫu và hy vọng Việt Nam đột phá để họ đi theo. Bây giờ họ đã thay đổi còn chúng ta co lại. Việt Nam phải mất thêm bao nhiêu thập niên nữa?!”. Thực tế phát triển các quốc gia đó sau ba thập niên chuyển đổi cho thấy tầm nhìn của ông là chính xác!

Từ khi rời Pháp, tôi chỉ chuyên tâm vào ngành chuyên môn hẹp của mình và không gặp lại ông. Năm 1993, từ Vancouver, được đọc bài góp ý của ông tại Hội nghị Ủy ban Trung ương Mặt trận Tổ quốc Việt Nam về việc sửa đổi Hiến pháp 1992, tôi càng cảm phục ông hơn. Trực tiếp, rõ ràng, bao quát mà cũng rất tập trung, các đề tài triết học về xã hội và chủ nghĩa phức tạp được ông nối kết với thực tế cụ thể của đời sống xã hội trước mắt bằng một cách trình bày mạch lạc, khoa học, thuyết phục. Phan Đình Diệu trước sau vẫn là một con người nhất quán!

Cho tới bây giờ, hơn hai mươi lăm năm sau, các ý kiến ông đóng góp vẫn còn rất giá trị, và càng ngày thực tiễn càng cho thấy chúng phù hợp với các giá trị của thế giới hướng về văn minh. Cho dù các ý kiến ông nêu không thuyết phục được hệ thống chính trị Việt Nam thời đó, chúng cho tôi thêm chứng cớ rằng trong xã hội Việt Nam không thiếu những tàng long ngọa hổ. Khi cuộc sống thực tế vấp vào những những tảng đá nằm dưới đáy con đường phát triển, thì những ý kiến như Phan Đình Diệu và những người khác đã chân thành đóng góp có thể sẽ có vai trò điểm tựa để dân tộc vươn lên.

Không chỉ dũng cảm trong suy nghĩ, chấp nhận cái mới, sẵn sàng loại bỏ quan điểm nào của mình đã chậm tiến, ông còn dũng cảm lên tiếng về điều ông nhận thức và tâm đắc, cho dù khác hẳn ý của những người hay thế lực cầm đại quyền. Đúng là một con người trí thức trung thực.

Con người đó rất sẵn sàng nhận trách nhiệm gánh vác việc xã hội, đất nước. Với tôi, ông có tầm vóc một nhà lãnh đạo đáng kính. Tiếc thay đất nước đã không dùng được tài năng và đức độ đó. Không cần ngó tới những vị trí cao hơn, nhìn tư cách, tác phong, tầm nhìn của ông mà tiếc biết bao cho ngành giáo dục nước nhà!

Giáo sư Phan Đình Diệu đã ra đi. Xin tiễn đưa anh với tấm lòng thương mến và kính phục của một người đang phải học theo anh rất nhiều về khía cạnh dũng cảm tri thức và tinh thần dấn thân.

Lê Học Lãnh Vân